\section{Method}
Suppose we have a time series in $\mathbb{R}^d$ over a period of discrete times $t_1,\cdots,t_T$, 
where for each time stamp we have $n_i$ observations. 
We assume the time series is observed from a Gaussian multivariate random vector $X\sim\mathcal{N}(0,\Sigma(t))$. 
To infer the dynamic network from the time series, 
we aim at estimating the sparse inverse covariance matrices $\Theta=\{\Theta_1,\cdots,\Theta_T\}$ from observations. 
This is achievable by minimizing the negative log-likelihood function with a temporal constraint and a sparsity penalty.
Hallac et al. define the loss function by
% \begin{equation*}
%     \mathcal{L} = \sum_{i=1}^{T}\left(-l(\Theta_i) + \lambda\|\Theta_i\|_{\text{od}, 1}\right) + 
%     \beta\cdot\left\{\begin{aligned}
%         &\sum_{i=2}^{T}\psi(\Theta_i-\Theta_{i-1}) & \text{Synchronous Observations}\\
%         &\sum_{i=2}^{T}\Delta_i\psi(\frac{\Theta_i-\Theta_{i-1}}{\Delta_i}) & \text{Asynchronous Observations}
%     \end{aligned}\right.,
% \end{equation*}
\begin{equation}
    \label{eq: ch1_TVGLloss}
    \mathcal{L}_{\text{TVGL}} = \sum_{i=1}^{T}\left(-l(\Theta_i) + \lambda\|\Theta_i\|_{\text{od}, 1}\right) + \beta\sum_{i=2}^{T}\psi(\Theta_i-\Theta_{i-1}), 
\end{equation}
where $\lambda$, $\beta$ are hyperparameters, $\|\cdot\|_{\text{od}, 1}$ is off-diagonal 1-norm, 
and $\psi(\cdot)$ is a convex penalty function, minimized at $\psi(0)$, 
encouraging similarity between adjacent sparse inverse covariance matrices. 

\subsection{Loss function of msTVGL}
Denote the sample covariance matrix by $S$. 
For each $t_i$, the negative log-likelihood function is given by 
\begin{align}
    \label{eq: ch1_loss}
    -l(\Theta_i) &= -\log\left((2\pi)^{-\frac{n_id}{2}}|\Theta_i|^{\frac{n_i}{2}}\exp\left\{-\frac{1}{2}\sum_{j=1}^{n_i}x_j^T\Theta_ix_j\right\}\right) \notag \\
               &= -\log((2\pi)^{-\frac{n_id}{2}}) - \frac{n_i}{2}\log(|\Theta_i|) +\frac{n_i}{2}\sum_{j=1}^{n_i}\text{Tr}(\frac{1}{n_i}x_j^T\Theta_ix_j) \notag \\
               &= -\log((2\pi)^{-\frac{n_id}{2}}) - \frac{n_i}{2}\log(|\Theta_i|) +\frac{n_i}{2}\text{Tr}(\frac{1}{n_i}\sum_{j=1}^{n_i}x_jx_j^T\Theta_i) \notag \\
               &= -\log((2\pi)^{-\frac{n_id}{2}}) - \frac{n_i}{2}\log(|\Theta_i|) +\frac{n_i}{2}\text{Tr}(S_i\Theta_i). 
\end{align}
Ignoring the constant and a scale, we have the loss used in paper of Hallac et al.\cite{hallac2017networkinferencetimevaryinggraphical}.
\begin{equation}
    \label{eq: ch1_loss_paper}
    -l(\Theta_i) = - n_i\log(|\Theta_i|) +n_i\text{Tr}(S_i\Theta_i). 
\end{equation}
Minimizing negative log-likelihood function encourages $\Theta_i=S_i^{-1}$. 

However, if for each time stamp there are multiple subjects contributing total observations, 
the negative log-likelihood function may not accurate anymore. 
Consider that we test on a certain type of cell in blood samples from multiple subjects. 
In this situation, cells from different subjects compose the set of observations. 
It is feasible to use the combined dataset to estimate the precision matrix but the heterogeneity between each subject isn't equally emphasized by the imbalanced sample sizes. 
In addition, if we only observe several positive definite matrices for each time stamp, 
we can't derive both sample size and the overall "sample covariance matrix". 
Hence, we modify the negative log-likelihood part of the loss function (\ref{eq: ch1_loss}) by 
using sample covariance matrices separately from different subjects while keeping a common precision matrix at each time stamp $t_i$,  
as presented as follows. 
\begin{align}
    \label{eq: ch1_loss_modify1}
    -l(\Theta_i) &= -\log((2\pi)^{-\frac{n_id}{2}}) - \frac{n_i}{2}\log(|\Theta_i|) +\frac{n_i}{2}\text{Tr}(S_i\Theta_i) \notag \\
                 &= \sum_{k=1}^{N_i}-\log((2\pi)^{-\frac{n_{ik}d}{2}})-\frac{n_{ik}}{2}\log(|\Theta_i|)+\frac{n_{ik}}{2}\text{Tr}(S_{ik}\Theta_i),
\end{align}
where there are $N_i$ subjects and $\sum_{k=1}^{N_i}n_{ik}=n_i$.

To reduce the influence of imbalanced sample sizes, 
we further normalize (\ref{eq: ch1_loss_modify1}) by removing the sample sizes $n_{i1},\cdots,n_{iN_i}$. 
This produces the final log-likelihood loss suitable for the multi-subject case (up to a constant and scale): 
\begin{align}
    \label{eq: ch1_loss_modify1_normalized}
    -l(\Theta_i) &= \sum_{k=1}^{N_i}-\log(|\Theta_i|)+\text{Tr}(S_{ik}\Theta_i) \notag \\
                 &= -N_i\log(|\Theta_i|)+N_i\text{Tr}(\bar{S_i}\Theta_i), 
\end{align}
where $\bar{S_i}$ is the mean of sample covariances of all subjects $\frac{1}{N_i}\sum_{k=1}^{N_i}S_{ik}$. 
Optimizing (\ref{eq: ch1_loss_modify1_normalized}) produces a precision matrix close to $\bar{S_i}^{-1}$.

The spatial penalty function is given by the off-diagonal 1-norm of each precision matrix, 
controlled by a penalty parameter $\lambda$, 
which can be written explicitly by $\lambda\|\Theta_i\|_{\text{od},1}=\lambda\sum_{j,k:j\neq k}|\Theta_{j,k}|$. 
It encourages the sparsity of the precision matrix.

For temporal penalty function $\psi$, 
Hallac et al.\cite{hallac2017networkinferencetimevaryinggraphical} proposed 5 choices, 
which are element-wise $l_1$ penalty, group lasso $l_2$ penalty, Laplacian penalty, 
$l_{\infty}$ penalty and row-column overlap penalty. 
Here element-wise $l_1$ penalty $\psi(X)=\sum_{i,j}|X_{i,j}|$ 
encourages the networks at adjacent time points to remain consistent on the majority of the edges, 
while only changing on a few edges. 
Group lasso $l_2$ penalty $\psi(X)=\sum_{j}\|[X]_j\|_2$ 
leads to a complete reconfiguration of the entire network at a few specific time points, 
when maintaining structural stability at the rest time points.
Laplacian penalty $\psi(X)=\sum_{i,j}X_{i,j}^2$ is suitable for smooth time-varying dynamic 
by allowing slight difference between adjacent time stamps and heavily punishing prominent differences. 
$l_{\infty}$ penalty $\psi(X)=\sum_{j}(\max_i|X_{i,j}|)$ is designed for scenarios where change happens only by clusters
because within the largest variation range of all the edges connected to a certain node, 
all edges in this cluster can freely vary without incurring additional penalties.
row-column overlap penalty $\psi(X)=\min_{V:V+V^T=X}\sum_{j}\|[V]_j\|_2$ is suitable for situations of repositioning of a single node to a new set of neighbors at each time stamp.
As we especially focus on research background of time-varying trend of cell-cell communication and brain region interaction, 
we discuss the most biologically meaningful temporal penalty types.

Usually, the time-varying dynamic of cell-cell interaction networks exhibits phased changes. 
Within each phase, the changes are smooth and continuous. 
Between different phases, there are sudden changes in edges and possibly network reorganization.
The reason is that those networks reflect the overall influence from 
the gradual establishment of cell communication patterns during development like immune system maturation during adolescence, 
the slow attenuation of communication intensity or the compensatory enhancement of specific pathways during aging, 
as well as the sudden activation or inhibition of certain signaling axes during disease progression. 
The change of brain structural connectivity network exhibits either a gradual smooth transition or sudden network reconfiguration.
The biological progression includes a continuous but gradually slowing down process of myelination and pruning of fiber bundles (with smooth and continuous changes) from childhood to adulthood;
possible relatively rapid global network reorganization before and after key developmental milestones (such as the maturation of language areas and prefrontal cortex functions).  
In special cases, the brain may suffer from a global connectivity disruptions in the course of neurodegenerative diseases (such as Alzheimer's disease), manifested as phased disintegration of the network.

Based on the aforementioned change characteristics of both types of networks, 
this paper first of all suggests Laplacian penalty functions because it aligns with the smooth and gradual change of the two types of networks, 
especially in studies with densely sampled time points.
Moreover, the element-wise $l_1$ penalty is adopted considering that in both types of networks, 
there are "switch-like" change events between a few key edges. 
For instance, the activation of a specific signaling pathway leads to a significant increase in the communication probability between certain cell types, 
or the sudden myelination of white matter tracts results in enhanced connectivity between certain brain regions.
At last, group lasso $l_2$ penalty works in the study since it is able to sniff out the sudden transition of cell-cell interaction pattern during certain stage in development, 
as well as to catch the large difference between networks when the gap between time points are huge, 
for example in studies across very rough time labels like "neonate", "adolescent", "adult" and "the elder". 
For the rest two penalty types $l_{\infty}$ and row-column overlap, 
since the coordinated changes of edges within a node cluster 
and the behavior of a single node reconnecting all edges at a single time point are extremely rare in cellular interaction networks and brain structural connection networks, 
and may even lack clear biological mechanism explanations, 
we do not intend to adopt them.

The above discussion on the negative loss function, spatial penalty and temporal penalty gives the final loss function of msTVGL,
as shown by (\ref{eq: ch1_lossmsTVGL})
\begin{equation}
    \label{eq: ch1_lossmsTVGL}
    \mathcal{L}_{\text{msTVGL}} = \sum_{i=1}^{T}(-N_i\log(|\Theta_i|)+N_i\text{Tr}(\bar{S_i}\Theta_i) + \lambda\|\Theta_i\|_{\text{od},1}) +
    \beta\sum_{i=2}^{T}\psi(\Theta_i-\Theta_{i-1}),
\end{equation}
where the subjects for all time points are $N_1,\cdots,N_T$ respectively, 
and temporal penalty $\psi$ is chosen from $\psi(\Theta_i-\Theta_{i-1})=\left\{\begin{aligned}
        &\sum_{j,k}|(\Theta_i-\Theta_{i-1})_{j,k}| & (\text{element-wise } l_1) \\
        &\sum_{j}\|[\Theta_i-\Theta_{i-1}]_j\|_2 & (\text{group lasso } l_2) \\
        &\sum_{j,k}(\Theta_i-\Theta_{i-1})^2_{j,k} & (\text{Laplacian})
    \end{aligned}\right..$

\subsection{ADMM based optimizing algorithm}
Hallac et al.\cite{hallac2017networkinferencetimevaryinggraphical} proposed an ADMM based algorithm to minimize the loss function (\ref{eq: ch1_TVGLloss}) with proximal operators. 

The proximal operator\cite{Parikh2014Proximal} of a convex function $f:\mathbb{R}^{m\times n}\mapsto\mathbb{R}\bigcup\{+\infty\}$ and a parameter $\eta\in\mathbb{R}$ is defined by 
\begin{equation*}
    \operatorname{prox}_{\eta f}(\mu):= \arg\min_{x\in\mathbb{R}^{m\times n}}(f(x)+\frac{1}{2\eta}\|x-\mu\|^2_F), \forall \mu\in\mathbb{R}^{m\times n}.
\end{equation*}
It's clear that $\operatorname{prox}_{\eta f}:\mathbb{R}^{m\times n}\mapsto\mathbb{R}$ is a minimization task with the trade-off between a function and a Frobenius distance. 
ADMM is an optimizing framework designed for separable convex optimization problems with linear constraints\cite{Boyd2011distributed}, 
i.e. optimization problems with a form of 
\begin{equation*}
    \min_{x,z} f(x)+g(z) \text{ s.t. } Ax+Bz=c, \text{ for convex }f,g.
\end{equation*}
It introduces a dual variable $u$ to construct an augmented Lagrangian function
\begin{equation*}
    L_{\rho}(x,z,y)=f(x)+g(z)+\left\langle u,Ax+Bz-c\right\rangle_{F}+\frac{\rho}{2}\|Ax+Bz-c\|_F^2, 
\end{equation*}
where $\|\cdot\|_F$ is Frobenius norm and $\rho>0$ is the ADMM penalty parameter. 
The variables are solved by an algorithm where at $k$-th step the updating rule is
\begin{align}
    x^{(k+1)} &:= \arg\min_x L_{\rho}(x,z^{(k)},u^{(k)}) \label{eq: ch1_ADMMupdating1} \\
    z^{(k+1)} &:= \arg\min_z L_{\rho}(x^{(k+1)},z,u^{(k)}) \label{eq: ch1_ADMMupdating2} \\
    u^{(k+1)} &:= u^k + \rho(Ax^{(k+1)}+Bz^{(k+1)}-c). \label{eq: ch1_ADMMupdating3}
\end{align}

To fit in the ADMM framework, Hallac et al. gave the equivalent optimization task of (\ref{eq: ch1_TVGLloss}) by 
introducing a consensus variable $Z=\{Z_0,Z_1,Z_2\}$ such that it forms three restrictive conditions $Z_0=\Theta$, 
$Z_1=\{\Theta_1,\cdots,\Theta_{T-1}\}$, $Z_2=\{\Theta_2,\cdots,\Theta_T\}$. 
Then they used the fact that the minimization tasks (\ref{eq: ch1_ADMMupdating1}) and (\ref{eq: ch1_ADMMupdating2}) are equivalent to several proximal operators, 
which have closed-forms, to derive the analytical form of the updating rule.
Since the loss function of msTVGL (\ref{eq: ch1_lossmsTVGL}) differs from (\ref{eq: ch1_TVGLloss}) only by the negative log-likelihood, 
we directly implement the updating rule of the consensus variable $Z$ and dual variable $U$ by Hallac et al., 
and derive the new closed-form updating rule of $\Theta$.

\subsubsection{Equivalent optimization task and augmented Lagrangian function}
After introducing the consensus variables $Z$, the optimization task of (\ref{eq: ch1_lossmsTVGL}) is equivalent to
\begin{align*}
    & \min_{\Theta, Z_0, Z_1, Z_2}\sum_{i=1}^{T}\left(-N_i\log(|\Theta_i|)+N_i\text{Tr}(\bar{S_i}\Theta_i) + \lambda\|Z_{i,0}\|_{\text{od}, 1}\right) + \beta\sum_{i=1}^{T-1}\psi(Z_{i,2}-Z_{i,1}) \\
    & \text{Subject to } Z_{i,0}=\Theta_i,\ (Z_{j,1}, Z_{j,2})=(\Theta_j,\Theta_{j+1}),\ 1\leq i\leq T,\ 1\leq j\leq T-1.
\end{align*}
The augmented Lagrangian function is given by 
\begin{align*}
    L_p(\Theta,Z,\tilde{U}) =& \sum_{i=1}^{T}\left(-N_i\log(|\Theta_i|)+N_i\text{Tr}(\bar{S_i}\Theta_i) + \lambda\|Z_{i,0}\|_{\text{od}, 1}\right) + \beta\sum_{i=1}^{T-1}\psi(Z_{i,2}-Z_{i,1}) \\
                                  &+ \frac{\rho}{2}\sum_{i=1}^{T}(\|\Theta_i-Z_{i,0}+\frac{1}{\rho}\tilde{U}_{i,0}\|_F^2-\|\frac{1}{\rho}\tilde{U}_{i,0}\|_F^2) \\
                                  &+ \frac{\rho}{2}\sum_{i=1}^{T-1}(\|\Theta_{i}-Z_{i,1}+\frac{1}{\rho}\tilde{U}_{i,1}\|_F^2-\|\frac{1}{\rho}\tilde{U}_{i,1}\|_F^2+\|\Theta_{i+1}-Z_{i,2}+\frac{1}{\rho}\tilde{U}_{i,2}\|_F^2-\|\frac{1}{\rho}\tilde{U}_{i,2}\|_F^2), 
\end{align*}
where $\tilde{U}=(\underbrace{(\tilde{U}_{1,0},\cdots,\tilde{U}_{T,0})}_{\tilde{U}_0},\underbrace{(\tilde{U}_{1,1},\cdots,\tilde{U}_{T-1,1})}_{\tilde{U}_1},\underbrace{(\tilde{U}_{1,2},\cdots,\tilde{U}_{T-1,2})}_{\tilde{U}_2})$
is the dual variable.
For simplicity, we denote $\frac{1}{\rho}\tilde{U}$ by $U$ and update it instead of $\tilde{U}$ in the algorithm.

\subsubsection{Updating $\Theta$}

Suppose at $k+1$-th step we have $Z^{(k)}$, $U^{(k)}$ and use them to update $\Theta^{(k+1)}$.
For simplicity, we denote $Z_{i,j}-U_{i,j}$ by $B_{i,j}$.

\textbf{i) For $i=2,\cdots,T-1$:} 
\begin{align*}
    \Theta_i^{(k+1)}&=\arg\min_{\Theta_i}-N_i\log(|\Theta_i|)+N_i\text{Tr}(\bar{S_i}\Theta_i)+\frac{\rho}{2}\|\Theta_i-B_{i,0}^{(k)}\|_F^2+\frac{\rho}{2}\|\Theta_i-B_{i,1}^{(k)}\|_F^2+\frac{\rho}{2}\|\Theta_{i}-B_{i-1,2}^{(k)}\|_F^2 \\
                    &=\arg\min_{\Theta_i}-N_i\log(|\Theta_i|)+N_i\text{Tr}(\bar{S_i}\Theta_i)+\frac{3\rho}{2}\left(\|\Theta_i\|_F^2-2\left\langle \Theta_i,\frac{1}{3}(B_{i,0}^{(k)}+B_{i,1}^{(k)}+B_{i-1,2}^{(k)})\right\rangle_F\right) + \underbrace{C(B_{i,0}^{(k)},B_{i,1}^{(k)},B_{i-1,2}^{(k)})}_{\text{Irrelevant to }\Theta_i} \\
                    &=\arg\min_{\Theta_i}-N_i\log(|\Theta_i|)+N_i\text{Tr}(\bar{S_i}\Theta_i)+\frac{1}{2\frac{1}{3\rho}}\|\Theta_i-A_i^{(k)}\|_F^2 + \underbrace{C'(B_{i,0}^{(k)},B_{i,1}^{(k)},B_{i-1,2}^{(k)})}_{\text{Irrelevant to }\Theta_i}, 
\end{align*}
where $A_i^{(k)}:=\frac{1}{3}(B_{i,0}^{(k)}+B_{i,1}^{(k)}+B_{i-1,2}^{(k)})$.
Hence, $\Theta_i^{(k+1)}=\operatorname{prox}_{\frac{1}{3\rho}(-N_i\log|\cdot|+N_i\text{Tr}(\bar{S}_i\cdot))}(A_i^{(k)})$, 
which has a closed-form expression. 
In fact, by setting gradient with respect to $\Theta_i$ to 0, we have 
\begin{align}
    \label{eq: ch1_Thetaupdate1}
    0 &= -N_i\Theta_i^{-1}+N_i\bar{S}_i+3\rho(\Theta_i-A_i^{(k)}) \notag \\
   \Longleftrightarrow -N_i\Theta_i^{-1} +3\rho\Theta_i &= 3\rho A_i^{(k)}-N_i\bar{S}_i \notag \\
   \Longleftrightarrow 0 &= \Theta_i^2 + (\frac{N_i}{3\rho}\bar{S}_i-A_i^{(k)})\Theta_i - \frac{N_i}{3\rho}I.
\end{align}
Denote the eigenvalue decomposition for $\frac{N_i}{3\rho}\bar{S}_i-A_i^{(k)}$ by $Q_iD_iQ_i^T$, 
where $Q_i$ is an orthonormal matrix of eigenvectors and $D_i$ is a diagonal matrix of eigenvalues. 
Then (\ref{eq: ch1_Thetaupdate1}) equals 
\begin{align}
    \label{eq: ch1Thetaupdate2}
    0 &= \Theta_i^2 + Q_iD_iQ_i^T\Theta_i - \frac{N_i}{3\rho}I \notag \\
    Q_i(\frac{1}{4}D_i^2+\frac{N_i}{3\rho}I)Q_i^T &= (\Theta_i+\frac{1}{2}Q_iD_iQ_i^T)^2.
\end{align}
By the positive definiteness of $\Theta_i$, 
we have $\Theta_i^{(k+1)}:=Q_i(-\frac{1}{2}D_i+\sqrt{\frac{1}{4}D_i^2+\frac{N_i}{3\rho}})Q_i^T$ by (\ref{eq: ch1Thetaupdate2}).

\vspace{1cm}
\textbf{ii) For $i=1$:} 
\begin{align*}
    \Theta_1^{(k+1)}&=\arg\min_{\Theta_1}-N_1\log(|\Theta_1|)+N_1\text{Tr}(\bar{S_1}\Theta_1)+\frac{\rho}{2}\|\Theta_1-B_{1,0}^{(k)}\|_F^2+\frac{\rho}{2}\|\Theta_1-B_{1,1}^{(k)}\|_F^2 \\
                    &=\arg\min_{\Theta_1}-N_1\log(|\Theta_1|)+N_1\text{Tr}(\bar{S_1}\Theta_1)+\frac{1}{2\frac{1}{2\rho}}\|\Theta_1-A_1^{(k)}\|_F^2 + \underbrace{C(B_{1,0}^{(k)},B_{1,1}^{(k)})}_{\text{Irrelevant to }\Theta_1}, 
\end{align*}
where $A_1^{(k)}:=\frac{1}{2}(B_{1,0}^{(k)}+B_{1,1}^{(k)})$. 
Hence, $\Theta_1^{k+1}=\operatorname{prox}_{\frac{1}{2\rho}(-N_1\log|\cdot|+N_1\text{Tr}(\bar{S}_1\cdot))}(A_1^{(k)})$, 
which has a closed-form expression $\Theta_1^{k+1}:=Q_1(-\frac{1}{2}D_1+\sqrt{\frac{1}{4}D_1^2+\frac{N_1}{2\rho}})Q_1^T$, 
with eigenvalue decomposition $\frac{N_1}{2\rho}\bar{S}_1-A_1^{(k)}:=Q_1D_1Q_1^T$.

\vspace{1cm}
\textbf{iii) For $i=T$:}
\begin{align*}
    \Theta_T^{(k+1)}&=\arg\min_{\Theta_T}-N_T\log(|\Theta_T|)+N_T\text{Tr}(\bar{S_T}\Theta_T)+\frac{\rho}{2}\|\Theta_T-B_{T,0}^{(k)}\|_F^2+\frac{\rho}{2}\|\Theta_T-B_{T-1,2}^{(k)}\|_F^2 \\
                    &=\arg\min_{\Theta_T}-N_T\log(|\Theta_T|)+N_T\text{Tr}(\bar{S_T}\Theta_T)+\frac{1}{2\frac{1}{2\rho}}\|\Theta_T-A_T^{(k)}\|_F^2+\underbrace{C(B_{T,0}^{(k)},B_{T-1,2}^{(k)})}_{\text{Irrelevant to }\Theta_T},
\end{align*}
where $A_T^{(k)}:=\frac{1}{2}(B_{T,0}^{(k)}+B_{T-1,2}^{(k)})$. 
Hence, $\Theta_T^{k+1}=\operatorname{prox}_{\frac{1}{2\rho}(-N_T\log|\cdot|+N_T\text{Tr}(\bar{S}_T\cdot))}(A_T^{(k)})$, 
which has a closed-form expression $\Theta_T^{k+1}:=Q_T(-\frac{1}{2}D_T+\sqrt{\frac{1}{4}D_T^2+\frac{N_T}{2\rho}})Q_T^T$, 
with eigenvalue decomposition $\frac{N_T}{2\rho}\bar{S}_T-A_T^{(k)}:=Q_TD_TQ_T^T$.

\subsubsection{Updating $Z$}
We implement directly the updating step of $Z$ by Hallac et al.\cite{hallac2017networkinferencetimevaryinggraphical}. 
To ensure the completeness of the content, 
we provide the derivations of the parts that were omitted and display the results that have already been shown in the original paper.
Suppose at $k+1$-th step we have $\Theta^{(k+1)}$, $U^{(k)}$ and use them to update $Z^{(k+1)}$. 

\textbf{Part i) Updating $Z_0$:}
\begin{align*}
    Z_{i,0}^{(k+1)}&=\arg\min_{Z_{i,0}}\lambda\|Z_{i,0}\|_{\text{od,1}}+\frac{\rho}{2}\|Z_{i,0}-\Theta_i^{(k+1)}-U_{i,0}^{(k)}\|_F^2 \\
                   &=\operatorname{prox}_{\frac{\lambda}{\rho}\|\cdot\|_{\text{od,1}}}(\Theta_i^{(k+1)}+U_{i,0}^{(k)}).
\end{align*}
To get the explicit expression of $\operatorname{prox}_{\frac{\lambda}{\rho}\|\cdot\|_{\text{od,1}}}(\Theta_i^{(k+1)}+U_{i,0}^{(k)})$, 
we derive by gradient. 
In fact, setting the gradient with respect to $Z_{i,0}$ to 0, 
we have 
\begin{align*}
    0&=\nabla_{Z_{i,0}}\|Z_{i,0}\|_{\text{od,1}} + \frac{\rho}{\lambda}(Z_{i,0}-(\Theta_i^{(k+1)}+U_{i,0}^{(k)})) \\
    \Longleftrightarrow Z_{i,0}^{(k+1)}+\frac{\lambda}{\rho}\nabla_{Z_{i,0}}\|Z_{i,0}^{(k+1)}\|_{\text{od,1}} &= \Theta_i^{(k+1)}+U_{i,0}^{(k)}.
\end{align*}
For entry $Z_{i,0}^{(k+1)}(j,l)$, $j\neq l$, we denote the sign of it by $\text{sgn}(Z_{i,0}^{(k+1)}(j,l))$. 
Then 
\begin{equation*}
    \left\{\begin{aligned}
        \text{sgn}(Z_{i,0}^{(k+1)}(j,l))=-1&:\ Z_{i,0}^{(k+1)}(j,l)=\underbrace{\frac{\lambda}{\rho}}_{>0}+(\Theta_i^{(k+1)}+U_{i,0}^{(k)})(j,l)<0 \\
        \text{sgn}(Z_{i,0}^{(k+1)}(j,l))= 1&:\ Z_{i,0}^{(k+1)}(j,l)=\underbrace{-\frac{\lambda}{\rho}}_{<0}+(\Theta_i^{(k+1)}+U_{i,0}^{(k)})(j,l)>0
    \end{aligned}\right.\Longrightarrow
    \left\{\begin{aligned}
        \text{sgn}((\Theta_i^{(k+1)}+U_{i,0}^{(k)})(j,l)) &= \text{sgn}(Z_{i,0}^{(k+1)}(j,l)) \\
        \bigg|(\Theta_i^{(k+1)}+U_{i,0}^{(k)})(j,l)\bigg| &> \frac{\lambda}{\rho}
    \end{aligned}\right.,
\end{equation*}
and by the definition of subgradient\footnote{The subgradient of $f:\mathbb{R}^n$ is $\partial f:=\{g:f(y)\geq f(x)+g^T(y-x),\forall y\in\mathbb{R}^n\}$.}
\begin{equation*}
    Z_{i,0}^{(k+1)}(j,l)=0:\ \nabla_{Z_{i,0}^{(k+1)}}\|Z_{i,0}\|_{\text{od,1}}=\frac{\rho}{\lambda}(\Theta_i^{(k+1)}+U_{i,0}^{(k)})\in[-1,1]\Longrightarrow(\Theta_i^{(k+1)}+U_{i,0}^{(k)})(j,l)\in[-\frac{\lambda}{\rho},\frac{\lambda}{\rho}].
\end{equation*}
Therefore, the closed-form expression of $Z_{i,0}^{(k+1)}$ is given by 
\begin{equation*}
    Z_{i,0}^{(k+1)}(j,l)=\left\{\begin{aligned}
        0,\ &\bigg|(\Theta_i^{(k+1)}+U_{i,0}^{(k)})(j,l)\bigg|\leq\frac{\lambda}{\rho},\ j\neq l \\
        \text{sgn}((\Theta_i^{(k+1)}+U_{i,0}^{(k)})(j,l))(|(\Theta_i^{(k+1)}+U_{i,0}^{(k)})(j,l)|-\frac{\lambda}{\rho}),\ &\bigg|(\Theta_i^{(k+1)}+U_{i,0}^{(k)})(j,l)\bigg|>\frac{\lambda}{\rho},\ j\neq l \\
        (\Theta_i^{(k+1)}+U_{i,0}^{(k)})(j,l),\ j=l
    \end{aligned}\right..
\end{equation*}

\vspace{1cm}
\textbf{Part ii) Jointly updating $Z_1,Z_2$:}
\begin{align*}
    Z_{i,1}^{(k+1)}, Z_{i,2}^{(K+1)}&=\arg\min_{Z_{i,1}, Z_{i,2}}\beta\psi(Z_{i,2}-Z_{i,1})+\frac{\rho}{2}\|Z_{i,1}-\Theta_i^{(k+1)}-U_{i,1}^{(k)}\|_F^2+\frac{\rho}{2}\|Z_{i,2}-\Theta_{i+1}^{(k+1)}-U_{i+1,2}^{(k)}\|_F^2 \\
                                    &=\arg\min_{Z_{i,1}, Z_{i,2}}\beta\psi([-I\ I]\begin{bmatrix}Z_{i,1}\\Z_{i,2}\end{bmatrix})+\frac{\rho}{2}\bigg\|\begin{bmatrix}Z_{i,1}-\Theta_i^{(k+1)}-U_{i,1}^{(k)}\\Z_{i,2}-\Theta_{i+1}^{(k+1)}-U_{i+1,2}^{(k)}\end{bmatrix}\bigg\|_F^2 \\
    &=\operatorname{prox}_{\frac{\beta}{\rho}\psi([-I\ I]\cdot)}(\begin{bmatrix}\Theta_i^{(k+1)}+U_{i,1}^{(k)}\\\Theta_{i+1}^{(k+1)}+U_{i+1,2}^{(k)}\end{bmatrix}) \\
    \Longrightarrow \begin{bmatrix}Z_{i,1}^{(k+1)}\\ Z_{i,2}^{(k+1)}\end{bmatrix} &=\frac{1}{2}\begin{bmatrix}\Theta_i^{(k+1)}+\Theta_{i+1}^{(k+1)}+U_{i,1}^{(k)}-U_{i+1,2}^{(k)}\\\Theta_i^{(k+1)}+\Theta_{i+1}^{(k+1)}+U_{i,1}^{(k)}-U_{i+1,2}^{(k)}\end{bmatrix}+\frac{1}{2}\begin{bmatrix}-\operatorname{prox}_{\frac{2\beta}{\rho}\psi}([\Theta_{i+1}^{(k+1)}-\Theta_i^{(k+1)}+U_{i+1,2}^{(k)}-U_{i,1}^{(k)}])\\\operatorname{prox}_{\frac{2\beta}{\rho}\psi}([\Theta_{i+1}^{(k+1)}-\Theta_i^{(k+1)}+U_{i+1,2}^{(k)}-U_{i,1}^{(k)}])\end{bmatrix},
\end{align*}
where the last step uses the property that if $f(x)=g(Cx)$, $CC^T=\frac{1}{\alpha}I$, 
then $\operatorname{prox}_f(x)=(I-\alpha C^TC)x+\alpha C^T\operatorname{prox}_{\frac{1}{\alpha}g}(x)$.

Hallac et al. used the property that if $f$ is block separable, i.e.$f(x)=\sum_{j}f_j(x_j)$, then $(\operatorname{prox}_f(y))_j=\operatorname{prox}_{f_j}(y_j)$, 
to give the closed-form expression of $E^{(k+1)}:=\operatorname{prox}_{\frac{2\beta}{\rho}\psi}([\Theta_{i+1}^{(k+1)}-\Theta_i^{(k+1)}+U_{i+1,2}^{(k)}-U_{i,1}^{(k)}])$. 
We denote $\Theta_{i+1}^{(k+1)}-\Theta_i^{(k+1)}+U_{i+1,2}^{(k)}-U_{i,1}^{(k)}$ by $A^{(k)}$ for simplicity.

\textbf{a)} When $\psi(\cdot):=\sum_{ij}|(\cdot)_{ij}|=\|\cdot\|_1$, the element-wise $l_1$ penalty, 
$E^{(k+1)}=\operatorname{prox}_{\frac{2\beta}{\rho}\|\cdot\|_1}(A^{(k)})$.
Following the derivation of the updating rule of $Z_{i,0}$, we have
$E^{(k+1)}_{ij}=\left\{\begin{aligned}
    0, &\bigg|A^{(k)}_{ij}\bigg|\leq\frac{2\beta}{\rho} \\
    \text{sgn}(A^{(k)}_{ij})(|A^{(k)}_{ij}|-\frac{2\beta}{\rho}), &\bigg|A^{(k)}_{ij}\bigg|>\frac{2\beta}{\rho}
\end{aligned}\right.$.

\textbf{b)} When $\psi(\cdot):=\sum_{j}\|[(\cdot)_j]\|_2$, the group lasso $l_2$ penalty,
$[E^{(k+1)}]_j=\operatorname{prox}_{\frac{2\beta}{\rho}\|(\cdot)_j\|_2}([A]_j)=\arg\min_{[E]_j}\|[E]_j\|_2+\frac{1}{2\frac{2\beta}{\rho}}\|[E]_j-[A]_j\|_F^2$. 
The setting gradient with respect to $[E]_j$ to 0 produces 
\begin{equation}
    \label{eq: ch1_gradientproxl2penalty}
    \frac{2\beta}{\rho}\nabla_{[E]_j}\|[E^{(k+1)}]_j\|_2+[E^{(k+1)}]_j=[A]_j. 
\end{equation}
If $[E]_j=0$, then by the definition of subgradient, $\nabla_{[E]_j}\|[E^{(k+1)}]_j\|_2=\{v:\|v\|_2\leq 1\}$, 
which means $[A]_j\in\{v:\|v\|_2\leq \frac{2\beta}{\rho}\}$.
If $[E]_j\neq 0$, then by (\ref{eq: ch1_gradientproxl2penalty})
\begin{align*}
    (\frac{2\beta}{\rho}\frac{1}{\|[E^{(k+1)}]_j\|_2}+1)[E^{(k+1)}]_j&=[A]_j \\
    \Longrightarrow\frac{2\beta}{\rho}+\|[E^{(k+1)}]_j\|_2&=\|[A]_j\|,
\end{align*}
which implies $[E^{(k+1)}]_j:=(1-\frac{2\beta}{\rho}\frac{1}{\|[A]_j\|_2})[A]_j$ is a solution.
Hence, $[E^{(k+1)}]_j=\left\{\begin{aligned}
    0, & \|[A]_j\|_2\leq\frac{2\beta}{\rho} \\
    (1-\frac{2\beta}{\rho}\frac{1}{\|[A]_j\|_2})[A]_j, &\|[A]_j\|_2>\frac{2\beta}{\rho}
\end{aligned}\right.$.

\textbf{c)} When $\psi(\cdot):=\sum_{ij}|(\cdot)|^2=\|(\cdot)\|^2_F$, the Laplacian penalty, 
we have $E^{(k+1)}_{ij}=\operatorname{prox}_{\frac{2\beta}{\rho}|\cdot|^2}(A^{(k)}_{ij})=\arg\min_{E_{ij}}|E_{ij}|^2+\frac{1}{2\frac{2\beta}{\rho}}(E_{ij}-A^{(k)}_{ij})^2$, 
which is just $E^{(k+1)}_{ij}:=\frac{1}{1+\frac{4\beta}{\rho}}A^{(k)}_{ij}$.

\subsubsection{Updating $U$} 
The updating rule of $U$ at $k+1$-th step follows from (\ref{eq: ch1_ADMMupdating3}), 
which is
\begin{equation*}
    U^{(k+1)}=\begin{bmatrix}U^{(k+1)}_0\\U^{(k+1)}_1\\U^{(k+1)}_2\end{bmatrix}:=
    \begin{bmatrix}U^{(k)}_0\\U^{(k)}_1\\U^{(k)}_2\end{bmatrix} + 
    \begin{bmatrix}
        \Theta^{(k+1)}-Z^{(k+1)}_0 \\
        \Theta^{(k+1)}\backslash\{\Theta^{(k+1)}_{T-1}\}-Z^{(k+1)}_1 \\
        \Theta^{(k+1)}\backslash\{\Theta^{(k+1)}_1\}-Z^{(k+1)}_2
    \end{bmatrix}.
\end{equation*}
Recall that $U=\frac{1}{\rho}\tilde{U}$, where $\tilde{U}$ is the dual variable, 
the ADMM penalty parameter $\rho$ is dropped on the both sides of the updating equation. 

\subsection{Hyperparameters}
\label{subsec: hyperparameters}
There are two positive hyperparameters controlling penalty strength. 
For sparsity penalty $\lambda\|\Theta\|_{\text{od}, 1}$, 
each off-diagonal element $x$ has gradient\footnote{Subgradient at $x=0$ is set to 0 but not $[-\lambda,\lambda]$.}
\begin{equation*}
    \left\{\begin{aligned}
        &\lambda, x>0 \\
        &0, x=0 \\
        -&\lambda, x<0
    \end{aligned}\right..
\end{equation*}
Since when optimizing convex penalty terms we expect zero gradients, 
the algorithm will push off-diagonal elements to zero hard as $\lambda$ gets large, 
and this result in a very sparse estimated inverse covariance matrix. 
When $\lambda$ is small, we expect $\Theta$ better matches the empirical data, 
but will lead to very dense networks, 
which are less interpretable and often overfit. 

This can be displayed using gradient. 
In fact, when we set temporal penalty to zero, 
then roughly, the gradient\footnote{Calculate by Fr\'echet derivative.} of the loss function (\ref{eq: ch1_loss_modify1_normalized}) plus the spatial penalty with respect to $\Theta_i$ is given by 
\begin{align*}
    \nabla_{\Theta_i}(-l(\Theta_i)+\lambda\|\Theta\|_{\text{od,1}}) &= \nabla_{\Theta_i}\left(-N_i\log(|\Theta_i|)+N_i\text{Tr}(\bar{S}_i\Theta_i) + \lambda\|\Theta_i\|_{\text{od}, 1}\right) \\
                                  &= -N_i\Theta_i^{-1} + N_i\bar{S}_i + (\pm\lambda)_{\text{od}},
\end{align*}
where $(\pm\lambda)_{\text{od}}$ represents a square matrix on $\mathbb{R}^{d\times d}$ with 0 diagonal and $\pm\lambda$ filling its off-diagonal. 
That 0 gradient is equivalent to $\Theta_i^{-1}=\bar{S}_i+\frac{1}{N_i}(\pm\lambda)_{\text{od}}$  
clarifies the influence of $\lambda$ on the deviation of estimated precision matrix from $\bar{S}_i^{-1}$.
In addition, the magnitude of $\lambda$ is associated with subject number $N_i$. 
If subject numbers are extremely imbalanced for all time points, 
a common sparsity penalty will cause over penalized precision matrices at time stamps with very few subjects. 
However, since for a real dataset the subject numbers usually don't vary drastically across all time points, 
this is not a problem that should be especially emphasized. 
In fact, the problem of imbalanced observations must be aware of in the original TVGL, 
because the gradient with respect to $\Theta_i$ of (\ref{eq: ch1_loss_paper}) plus the spatial penalty gives the relationship $\Theta_i^{-1}=S_i+\frac{1}{n_i}(\pm\lambda)_{\text{od}}$. 
Unlike subject numbers, the sample numbers tend to be extremely unequal for different time points, 
whereby over penalized precision matrices are produced. 
Hence, the loss function of msTVGL (\ref{eq: ch1_lossmsTVGL}) helps to alleviate this phenomena. 

For temporal penalty term, $\beta$ controls the transition between adjacent estimated inverse covariance matrices. 
Small $\beta$ will lead to fluctuating $\Theta$ between time stamps 
while large $\beta$ makes transition smooth. 
As $\beta\rightarrow\infty$, 
the problem becomes a static graphical lasso problem since a constant $\Theta$ is the solution.

\subsection{Hyperparameter selection}

The hyperparameter selection criterions such as AIC and BIC used by Hallac et al.\cite{hallac2017networkinferencetimevaryinggraphical}
are suggested based on the assumption of i.i.d. Gaussian observations. 
Since we extend the method to abstract situations of using symmetric positive definite matrices directly, 
we suggest an alternative hyperparameter selection method leveraging prior knowledge on the networks 
that the method learns. 
% For example, to generate an interpretable and biologically meaningful result, 
% the algorithm must assign the estimated precision matrices a rational sparsity to highlight significant edges and avoid learning too many noises from observations. 
% In addition, the temporal smoothness penalty must be given a proper strength to maintain consecution between adjacent time stamps. 
To achieve this, we introduce the average degree of nodes, Jaccard similarity of adjacent graphs and temporal deviation ratios as standards to restrict the two penalty coefficients. 

The average degree of nodes in a graph is defined by the mean number of edges connected to each node. 
Suppose there are $L$ edges in the graph and $p$ nodes, for an undirected graph, 
the average degree is given by 
\begin{equation*}
    \text{AD} = \frac{2L}{p}, 
\end{equation*}
where the coefficient 2 is due to each pair is checked twice by the edge. 
This measure represents the number of biologically meaningful edges in an interaction network. 
Usually, we set a threshold by prior knowledge and perceive entries with absolute values larger than this threshold in a precision matrix to be biologically 
meaningful edges. 
For different biologically backgrounds, the range of acceptable average degree varies, 
which also requires prior knowledge to determine.

The Jaccard similarity index measures the similarity between two sets. 
In the graph connectivity perspective, it represents the overlap rate of the two adjacent edge sets. 
Suppose $E_i$ and $E_{i+1}$ are two edge sets of the two graphs at adjacent time stamps $i$ and $i+1$, 
their Jaccard similarity index is given by 
\begin{equation*}
    J(i,i+1) = \frac{|E_i\bigcap E_{i+1}|}{|E_i\bigcup E_{i+1}|}.
\end{equation*} 
The Jaccard similarity index lies in $[0,1]$. 
When $J(i,i+1)$ is almost 1, the adjacent graphs are almost identical which implies a relatively strong temporal smoothness penalty. 
When $J(i,i+1)$ is close to 0, there are almost no shared edges between the two graphs. 
The range of Jaccard similarity varies a lot depending on the biological context, 
which need strong prior knowledge to set.

The temporal deviation ratio describes the sharp transitions between graphs. 
It is calculated by the average difference of two adjacent precision matrices in Frobenius norm. 
For precision matrices $\Theta_i$ and $\Theta_{i+1}$, the temporal deviation ratio at time shift $i\rightarrow i+1$ among $T$ time stamps is given by
\begin{equation*}
    \text{TDR}(i, i+1) = \frac{\|\Theta_{i+1}-\Theta_i\|_F}{\sum_{t=1}^{T-1}\|\Theta_{t+1}-\Theta_{t}\|_F}.
\end{equation*} 
We consider the time shift with a temporal deviation ratio above the $1.25$ times of the average a sharp transition. 

We choose the acceptable hyperparameter combinations by making sure
i) the average degree is within the required range which is featured by minimal and maximal acceptable percentages of node number; 
ii) the average Jaccard similarity lies in the satisfactory range $[\text{AD}_{l},\text{AD}_u]$;
iii) the temporal deviation ratio displays the least but at least 1 sharp transitions among all the attempts. 

The selection procedure is shown in Algorithm \ref{alg: ch1_hyperparametersel}. 

\begin{algorithm}
    \caption{Hyperparameter selection with warm start}
    \begin{algorithmic}
        \State 1) $S_{\lambda}\longleftarrow$ Searching space of $\lambda$.
        \State 2) With hyperparameter combination $(\lambda, \beta)$:
        \State \quad 2a) $\text{AD}_{\lambda,\beta}\longleftarrow$ Mean value of average degrees of all precision matrices at all time stamps.
        \State \quad 2b) $\text{TDR}_{\lambda,\beta}\longleftarrow$ Temporal deviation ratios of all precision matrices at all time stamps.
        \State \quad 2c) $\text{ST}_{\lambda,\beta}\longleftarrow$ Sharp transition numbers. 
        \State \quad 2d) $J_{\lambda,\beta}\longleftarrow$ Mean value of Jaccard similarities of all precision matrices at all time stamps.
        \State 3) $\text{AD}_{\text{min}}, \text{AD}_{\text{max}}\longleftarrow$ $\text{min}\%, \text{max}\%$ percentiles of node number. 
        \For{$\lambda, i\in \text{Enumerate}(S_{\lambda})$}
            \For{$\beta\in S_{\lambda}[0:i+1]$}
            \State a) Train model with $(\lambda, \beta)$.
            \State b) Compute $\text{AD}_{\lambda,\beta}$, $\text{TDR}_{\lambda,\beta}$, $J_{\lambda,\beta}$
            \EndFor
            \If{$\text{AD}_{\lambda^*,\beta^*}\in[\text{AD}_{\text{min}}, \text{AD}_{\text{max}}]$, $\text{ST}_{\lambda^*,\beta^*}=\arg\min_{\lambda,\beta}\{\text{ST}_{\lambda,\beta}>0\}$, $J_{\lambda^*,\beta^*}\in[\text{AD}_{l},\text{AD}_u]$}
                \State Accept $(\lambda^*, \beta^*)$
            \EndIf
        \EndFor
        \State \Return All accepted hyperparameters
    \end{algorithmic}
    \label{alg: ch1_hyperparametersel}
\end{algorithm}